\section{Classical Artin machinery and the scoped contribution}
\label{sec:classical}

\paragraph{Shift determinants.}
For a finite weighted shift of finite type, periodic-path traces and adjacency
determinants are two descriptions of the same zeta data.  The determinant and
rationality framework goes back at least to
\citet{BowenLanford1970}.  SD-C19 uses its smallest weighted instance: one base
vertex, parallel subset edges, and a finite group fiber.  The finite matrix is
not presented as a new definition of dynamical zeta.

\paragraph{Group extensions and twisted transfer operators.}
Representation-resolved dynamical \(L\)-functions for genuine group extensions
are classical.  \citet{ParryPollicott1986} develop Frobenius and Chebotarev
structure for Galois covers of Axiom-A flows;
\citet{AdachiSunada1987} establish twisted Perron--Frobenius and \(L\)-function
machinery; and Chapter 8 of \citet{ParryPollicott1990} gives a systematic
account of group extensions and their character factors.  Those works justify
the regular/isotypic decomposition used below.  We do not claim priority for
that mechanism.

\paragraph{Cohomology and periodic data.}
The periodic obstruction to a cocycle coboundary is rooted in Livšic theory
\citep{Livsic1972}.  Compact-group and matrix versions sharpen this principle
in hyperbolic settings \citep{ParryPollicott1997,Kalinin2011}.  The direction
needed here is elementary: a coboundary has identity product on every periodic
word, whereas the singleton fixed point of \eqref{1.1} has product \(a\ne e\).
This observation prevents a gauge-trivial construction from being counted as
fiber motion.

\paragraph{Graph and hypergraph covers.}
Artin factorization for finite graph coverings is developed by
\citet{StarkTerras1996,StarkTerras2000}.  The hypergraph extension of
\citet{EylerJun2024} is particularly close in vocabulary because its objects
also carry incidence and subset structure.  In those theories, local factors
come from covering prime cycles and their Frobenius elements.  SD-C19 instead
starts with a signed complete subset alphabet and a uniform degree cocycle.
The algebraic resemblance is real, but the orbit ledger is different.

\paragraph{A determinant is not a complete extension invariant.}
\citet{BoyleSchmieding2017} show that finite-group extensions of shifts of
finite type can share zeta or periodic data without being conjugate.  This is
an important claim boundary: the exact character determinants below do not
classify the cocycle or its symbolic extension.

\begin{table}[t]
\centering
\caption{Classical components and the contribution claimed for SD-C19.}
\label{tab:literature-boundary}
\small
\begin{tabularx}{\textwidth}{L{0.24\textwidth}
  >{\raggedright\arraybackslash}X
  >{\raggedright\arraybackslash}X}
\toprule
Component & Classical boundary & SD-C19 claim \\
\midrule
Finite shift determinant
& weighted adjacency and periodic traces
& signed tensor-subset specialization \\
Finite-group Artin blocks
& regular representation decomposes into irreducibles
& intrinsic degree cocycle repairs frozen-roof commutation \\
Cocycle cohomology
& periodic identity is necessary for a coboundary
& singleton fixed point proves genuine motion \\
Cover local factors
& Frobenius factors for graph/hypergraph prime cycles
& atom-local factor plus primitive-lift mismatch \\
Rigidity
& no general novelty claim
& one-letter natural operator-clean rules have cyclic power image \\
\bottomrule
\end{tabularx}
\end{table}

The contribution is therefore an exact synthesis plus a scoped theorem.  The
specific signed tensor-subset transfer admits a lawful parity cover, and the
same coefficient calculation that makes the cover clean proves that the
natural one-letter branch is only cyclic degree count.  Transition-dependent
cocycles remain outside this conclusion.
